\begin{figure}[t]
  \centering
  \begin{tikzpicture}[
    >=Latex,
    node distance=5mm and 7mm,
    source/.style={draw, rounded corners, fill=gray!12, align=center,
      font=\scriptsize, text width=1.45cm, minimum height=8mm},
    cell/.style={draw, rounded corners, fill=blue!12, align=center,
      font=\scriptsize, text width=2.0cm, minimum height=12mm},
    downstream/.style={draw, fill=green!12, align=center, font=\scriptsize,
      text width=1.4cm, minimum height=10mm},
    arrow/.style={->, semithick}
  ]
    \node[source] (a) {descriptor A\\cell index 0};
    \node[source, below=of a] (b) {descriptor B\\cell index 1};
    \node[cell, right=9mm of $(a)!0.5!(b)$] (cell) {one canonical cell\\same six scalars\\\texttt{CELL\_}$h$};
    \node[downstream, right=of cell] (down) {realization, KCs,\\items, Q-row};

    \draw[arrow] (a) -- (cell);
    \draw[arrow] (b) -- (cell);
    \draw[arrow] (cell) -- (down);
  \end{tikzpicture}
  \caption{Schematic deduplication with retained provenance.  Each incoming
  arrow is retained as a separate source edge even when two grammatical tuples
  share one content-derived canonical identifier.}
  \label{fig:provenance}
\end{figure}
